% =====================================================================
% STAGED DRAFT -- tier-1b two-asset HANK results paragraph.
% DO NOT \input THIS INTO THE PAPER YET.
%
% GATE: merge into green_deficits_price_level.tex ONLY when
%   output/tables/hank2_validation.txt shows
%       refinement (...): max rel. dev X.XXX -> PASS
% Both accuracy gates (oscillation + refinement) must read PASS. The
% oscillation gate already passes (scores 1, all four regimes). If the
% refinement pass FAILS, discard this draft and keep the "magnitudes NOT
% YET REPORTED" status in the paper.
%
% Numbers below are from the VERIFIED baseline run
% (output/tables/hank2_irfs_summary.txt, 2026-07-07): they become the
% reported numbers iff the refinement pass confirms them within 10%.
% If the refinement shifts any number, update it here from the refined
% grid before merging.
%
% NOTE: the refined grid (horizon/nb/na) is printed verbatim in the
% "refinement (TAYLOR, THORIZON ..., nb ..., na ...)" line of
% hank2_validation.txt. The default refined na is 36 (overridable via
% ACC_NA); set the "$n_a=..$" in the sentence below to whatever that line
% reports, and put the reported "max rel. dev" in place of [XX].
%
% WHERE IT GOES: replace the final bracket of the "A two-asset tier"
% paragraph -- currently
%   "[TIER-1b: IMPLEMENTED; all four regimes solve and pass the
%    oscillation diagnostic; refinement re-solve pending; magnitudes NOT
%    YET REPORTED.]"
% -- with the paragraph + table + figure below.
% =====================================================================

With both accuracy gates cleared (oscillation scores of one on every
regime; the refined re-solve --- horizon 500, illiquid grid $n_a=36$ ---
agreeing with the baseline within [XX]\%), the two-asset tier is
reportable, and it delivers the robustness the extension was built for.
The calibration clears the wage Phillips curve and the total-wealth market
at $\beta^{*}=0.971$, $\varphi^{*}=1.58$, with the portfolio-adjustment
cost held at the reference value $\chi_1=6.42$; grids are $n_e=3$,
$n_b=15$, $n_a=30$ (coarse, so magnitudes are indicative). Three findings
carry over from the one-asset tier essentially unchanged, and two are new.
\emph{First, the real green path is again rule-independent}: green capital
after forty quarters is $+0.165$ in \emph{every} regime (identical to
three decimals), and damages fall by $-0.0039$ throughout --- the monetary
rule moves nominal quantities, not the accumulation of abatement capital,
exactly as in the RANK and one-asset HANK tiers. \emph{Second, deficit
financing buys the debt path, not the inflation path}: the balanced-budget
comparator (\textsc{taylorbal}) accumulates real debt of only $+0.008$
after forty quarters against $+0.074$--$+0.077$ under the slow
debt-stabilizing rule --- a tenfold gap in the debt response --- while
impact inflation barely moves. \emph{Third, adding portfolio choice, a
capital stock, and sticky wages \underline{dampens} the impact inflation}:
under the Taylor rule it is $+0.04\%$ annualized (against $+0.44\%$ in the
one-asset tier), so the richer supply side absorbs more of the shock and
reinforces the paper's message that the inflationary margin is the
\emph{financing}, not the program. The two new objects are the point of
the tier. The \emph{liquid-bond return} $r^b$ falls on impact (e.g.\
$-0.0017$ under \textsc{weak}), mirroring inflation almost exactly: the
Fisher revaluation lands specifically on households' \emph{liquid} nominal
savings --- the $B/P$ margin the one-asset tier could not separate ---
rather than on total wealth. And the \emph{equity price} $p$ drops
$5$--$7\%$ on impact as the illiquid asset reprices against the program's
future debt and tax path. The two-asset tier thus confirms the one-asset
findings under a substantially richer household portfolio and production
block, and makes the liquid-safe-asset margin at the heart of the paper
directly observable. [TIER-1b LINEARIZED TWO-ASSET HANK IRF; run VERIFIED,
both accuracy gates passed; \texttt{hank2\_irfs\_summary.txt},
\texttt{hank2\_validation.txt}.]

\begin{table}[ht]
\centering
\caption{Tier-1b two-asset HANK: impulse responses to the quasi-permanent
green-investment shock ($\rho_g=0.98$, one percent of output) under four
monetary/financing regimes. Impact inflation and the liquid return are
per quarter; debt, green capital, and damages are at forty quarters.
Green capital is rule-independent; the balanced-budget regime
(\textsc{taylorbal}) isolates the deficit-financing component.
Generated by \texttt{run\_green\_hank2}.}
\label{tab:tier1b}
\begin{tabular}{lccccc}
\toprule
regime & $\pi$ impact (ann.) & $Y$ impact & $r^b$ impact & debt (40q) & $K_g$ (40q) \\
\midrule
\textsc{weak}       & $+0.67\%$ & $+0.0048$ & $-0.0017$ & $+0.074$ & $+0.165$ \\
\textsc{taylor}     & $+0.04\%$ & $+0.0052$ & $-0.0001$ & $+0.077$ & $+0.165$ \\
\textsc{greenaccom} & $+0.27\%$ & $+0.0054$ & $-0.0007$ & $+0.076$ & $+0.165$ \\
\textsc{taylorbal}  & $+0.03\%$ & $+0.0026$ & $-0.0001$ & $+0.008$ & $+0.165$ \\
\bottomrule
\end{tabular}
\end{table}

\pfigguard{PFig17_hank2_green_irfs.pdf}{%
\begin{figure}[ht]\centering
\includegraphics[width=0.95\textwidth]{PFig17_hank2_green_irfs.pdf}
\caption{Tier-1b two-asset HANK: linearized sequence-space impulse
responses (output, inflation, government debt, the liquid-bond return,
green capital, the equity price) to the quasi-permanent
deficit-financed green-investment shock under four regimes. The
liquid-return panel isolates the Fisher revaluation on households'
\emph{liquid} nominal savings; the equity-price panel is the illiquid
asset's repricing. Scope: linearized IRFs, not the nonlinear DTPL
price-level transition. Generated by \texttt{run\_green\_hank2}.}
\label{fig:pfig17}
\end{figure}}
